% ============================================================
%  Part 3: Scorpio 深度研究
% ============================================================

\section{Scorpio 深度研究}

\subsection{Scorpio 是什么}

Scorpio 不是传统 PCIe Switch，也不是 NVSwitch 的复制品。它是一种\textbf{新型的 Smart Fabric Switch}——基于 PCIe 6 物理层（64 GT/s PAM4）但可以运行多种协议：
标准 PCIe 6、CXL 3.x、或 memory semantic 的 platform-specific 协议。

\begin{table}[htb]
    \centering
    \caption{Scorpio 系列产品对比}
    \label{tab:scorpio}
    \begin{tabular}{lcc}
        \toprule
        \textbf{特性} & \textbf{Scorpio X-Series} & \textbf{Scorpio P-Series} \\
        \midrule
        定位 & Memory Semantic AI Fabric Switch & PCIe 6 Fabric Switch \\
        Lane Count & 320 lanes & 32–320 lanes \\
        协议 & Open \& Platform-Specific Memory Semantic & PCIe 6 / CXL 3.x \\
        目标应用 & GPU Scale-Up & 通用 I/O Fabric (GPU/NIC/NVMe) \\
        拓扑 & High-Radix, Single-Hop & 灵活 PCIe Fabric \\
        关键特性 & HW-Accelerated Collective Ops (Hypercast) & PCIe 6 Full Support \\
        出货状态 & 已出货至 Leading Hyperscalers & 已 Demo (H200 + ConnectX-8) \\
        \bottomrule
    \end{tabular}
\end{table}

Scorpio X-Series 的 320 lanes 是业界已知最大 radix 的 PCIe-based switch。
320 lanes 可以配置为 20 个 x16 端口、40 个 x8 端口、或 80 个 x4 端口。
Broadcom 的 PEX 系列公开产品最高约 144 lanes。高 radix 的核心优势是
\textbf{减少级联}：更多设备可以直接 single-hop 连接，无需多级交换级联，从而降低延迟。

\subsection{技术架构}

Scorpio 的技术分层：

\begin{enumerate}
    \item \textbf{物理层}：320$\times$ SerDes @ 64 GT/s PAM4（PCIe 6 PHY）
    \item \textbf{交换层}：High-Radix Crossbar Switch Fabric
    \item \textbf{协议引擎}：
        \begin{itemize}
            \item Standard PCIe 6 Protocol Engine（FLIT mode）
            \item Memory Semantic Protocol Engine（open + platform-specific）
        \end{itemize}
    \item \textbf{加速层}：HW-Accelerated Collective Operations (Hypercast)
    \item \textbf{管理层}：Fabric Manager via COSMOS Integration
\end{enumerate}

\textbf{关键技术参数（已确认）} \sourceT{2}：
\begin{itemize}
    \item 320 SerDes lanes @ 64 GT/s PAM4
    \item High-radix single-hop topology
    \item HW-accelerated collective ops (All-Reduce, All-to-All 等)
    \item Per-port telemetry via COSMOS
\end{itemize}

\textbf{未确认参数} \unverified{}：
\begin{itemize}
    \item 端到端转发延迟（未公开具体 ns 数）
    \item 功耗（未公开 W 数）
    \item Die size / 制程节点（未公开，推测 TSMC 7nm/5nm 级别）
\end{itemize}

\subsection{为什么 AI GPU 集群需要 Scorpio}

\textbf{场景 1：GPU Scale-Up Fabric}。多个 GPU 需要共享内存访问。NVSwitch 只支持 NVIDIA GPU。
异构集群（如混合 NVIDIA + AMD GPU 或 Google TPU + 其他加速器）需要开放替代方案。
自研 ASIC（如 Amazon Trainium、Google TPU）也需要 off-the-shelf 的 scale-up 互连方案。

\textbf{场景 2：GPU-to-NIC-to-Storage}。GPU 结果需快速写入 NVMe，或通过 NIC 发送到其他节点。
PCIe Switch 消除 CPU bottleneck——数据路径直接从 GPU 到 NIC，不经过 CPU。

\textbf{场景 3：MoE 模型通信优化}。Mixture-of-Experts 模型的每一层只有部分 expert 被激活
（例如 8 个中的 2 个），但 expert 分布在不同 GPU 上，需要频繁的 all-to-all 通信。
通信模式是``稀疏但大带宽''的——非常适合 fabric switch。
Astera Labs 官方博客（Caleb Shetland, Sr. Principal Firmware Architect）专门讨论了这个问题 \sourceT{2}。
Hypercast 技术在交换机内部对数据进行聚合（如 All-Reduce 的求和），
而非全部发到端点再计算，减少网络流量并降低延迟。
但具体的``tokens-per-watt''提升数据未公开 \unverified{}。

\textbf{场景 4：CXL 内存池}。GPU 通过 Scorpio 访问 CXL-attached DDR5 内存，
扩展有效内存容量（超越 HBM 限制）。关键 trade-off 是延迟 vs 容量。

\subsection{Scorpio vs NVSwitch 深度对比}

\begin{table}[htb]
    \centering
    \caption{Scorpio X-Series vs NVIDIA NVSwitch}
    \label{tab:scorpio-vs-nvswitch}
    \begin{tabular}{lcc}
        \toprule
        \textbf{对比维度} & \textbf{Scorpio X-Series} & \textbf{NVIDIA NVSwitch (H100/B200)} \\
        \midrule
        物理层 & PCIe 6 (64 GT/s PAM4) & NVLink 4/5 (100/200 GT/s) \\
        Per-GPU 带宽 & $\sim$128 GB/s (x16) & 900 GB/s (18 links $\times$ 50 GB/s) \\
        协议 & 开放 (PCIe/CXL) + Custom & 私有 (NVLink) \\
        GPU 兼容 & 任何 PCIe 设备 & 仅 NVIDIA GPU \\
        延迟 & 微秒级 (推测) & $<$1 微秒 \\
        Radix & 320 lanes (max) & 连接 8 GPU (DGX), 可级联 \\
        Collective Ops & HW-Accelerated (Hypercast) & SHARP (in-network compute) \\
        软件生态 & COSMOS & NVSwitch Fabric Manager \\
        \bottomrule
    \end{tabular}
\end{table}

Scorpio 在绝对带宽和延迟上无法与 NVSwitch 竞争（差距约 7$\times$ 带宽）。
其价值在于\textbf{开放性}、\textbf{异构兼容性}和\textbf{高 radix 灵活拓扑}。

NVLink/NVSwitch 带宽对比的量化理解：
\begin{itemize}
    \item NVSwitch (H100): 每个 GPU 有 18 个 NVLink 4 端口，每个 50 GB/s 双向 = 900 GB/s per GPU
    \item Scorpio via PCIe 6 x16: 64 GT/s $\times$ 16 lanes $\times$ 128/130 encoding = $\sim$128 GB/s per GPU
    \item 对于非 NVIDIA GPU（AMD MI300X PCIe、Intel Gaudi 3 PCIe），PCIe 是唯一的 scale-up 选项
    \item 即使是 H200，其 PCIe 侧也仍然是 x16 PCIe 5（64 GB/s）——在标准协议侧，Scorpio 匹配 GPU 的 PCIe 能力
\end{itemize}

\subsection{Scorpio vs Broadcom PCIe Switch}

\begin{table}[htb]
    \centering
    \caption{PCIe Fabric Switch 竞争对比}
    \label{tab:scorpio-vs-broadcom}
    \begin{tabular}{lcc}
        \toprule
        \textbf{维度} & \textbf{Scorpio (Astera)} & \textbf{PEX 系列 (Broadcom)} \\
        \midrule
        PCIe 世代 & PCIe 6 (前沿) & PCIe 5 为主, 6 roadmapped \\
        最大 lane 数 & 320 lanes & $\sim$144 lanes (公开) \\
        Memory Semantic & 支持 & 标准 PCIe only \\
        HW Collective Ops & 支持 (Hypercast) & 未知 \\
        AI 优化 & 从零为 AI 设计 & 通用数据中心 \\
        COSMOS 集成 & 原生 & 无等效平台 \\
        市场份额 & 新进入者 & 市场领导者 \sourceT{3} \\
        成熟度/稳定性 & 较新 & 极其成熟 \\
        \bottomrule
    \end{tabular}
\end{table}

Broadcom 的 PEX 系列（原 PLX Technology，2014 年被 Avago/Broadcom 收购）是行业 gold standard，
已经验证在数百万系统中，可靠性极高。Astera 的突破口在于 PCIe 6 是全新标准
（所有人都在新的起跑线上）、AI 集群的特殊需求、以及直接服务 hyperscaler 的能力。

\textbf{威胁评估}：对 Broadcom 的威胁程度为\textbf{中等}。
Scorpio 主要威胁的是 AI-specific 的增量市场，而非 Broadcom 的传统企业/存储 switch 业务。
但 AI market 的增速远超传统市场。

\subsection{Scorpio 的护城河与真实价值}

\begin{itemize}
    \item \textbf{SerDes 壁垒（最真实）}：320 lane 的 64 GT/s PAM4 SerDes 集成在一个芯片上是极大的模拟工程挑战。PAM4 的信噪比要求远高于 NRZ，crosstalk management 在 320 lane 密度下极其困难
    \item \textbf{Hyperscaler Qualification}：已通过 hyperscaler 的严格验证并出货。Qual 周期长达 12-24 个月，一旦通过更换成本高。但这只是时间壁垒，竞争对手有足够资金和时间也能通过
    \item \textbf{Memory Semantic Protocol}：可编程协议引擎的差异化是真实的，但如果竞争对手支持标准 CXL 3.x（也支持 memory semantic），功能会趋同
    \item \textbf{COSMOS 锁定}：统一管理的价值真实存在，但大型 hyperscaler 可以自己开发管理工具
\end{itemize}

\textbf{Scorpio 在 AI Infra 中的真实价值}：
目前是唯一基于 PCIe 并针对 AI 优化的大 radix fabric switch，
解决了异构 GPU 环境的 scale-up 互连问题，为 hyperscaler 提供了 NVIDIA 之外的战略性替代方案。
但其性能天花板（PCIe 带宽/延迟）限制了在最高性能场景的适用性。
